\section{Related Literature}
\label{sec:literature}

This paper sits at the intersection of four literatures.

\paragraph{Bid-coordination tests and bidder classification.}
\citet{bajari2003deciding} develop powerful tests for
collusion---exchangeability of bid functions and conditional
independence of bid residuals---but the researcher must first partition
bidders into ``suspected colluders'' and a ``competitive fringe.''
Subsequent work detects bid rotation
\citep{porter1993detection, porter1999ohio}, collusion in timber
auctions \citep{baldwin1997bidder}, bidder groups from co-bidding
patterns \citep{conley2016detecting}, and collusion dynamics in
long-lived cartels \citep{clark2021collusion,
schurter2020identification}. \citet{conley2016detecting} are the
closest precedent in data requirements, though they detect group
structure rather than flag individual firms. In each case the
partition is constructed \emph{ex post} or by assumption; the FL
classification provides a data-driven, \emph{ex ante} partition
grounded in a participation anomaly.

\paragraph{Proactive screens for bid rigging.}
\citet{imhof2019detecting} and \citet{imhof2018screening} train ML
classifiers on tender-level features and achieve high accuracy against
known cartels; see also \citet{abrantesmetz2006a},
\citet{harrington2008detecting}, \citet{chassang2022robust}, and
\citet{wallimann2023machine}. All operate at the \emph{tender} level
and require granular bid values. The FL screen operates at the
\emph{firm} level using only participation and outcome records,
enabling a two-stage workflow: FL flags triage firms, bid-level
analysis follows on flagged tenders.%
\footnote{Our Imhof comparison uses a CV median split, a coarser proxy
for the full implementation; see Section~\ref{sec:results_detection}.}

\paragraph{Cover bidding.}
Cover bidders have been documented in prosecuted cartels
\citep{porter1999ohio, pesendorfer2000study, asker2010leniency,
marshall2012economics} and in theoretical comparisons of auction
formats \citep{athey2011comparing}. In every case, identification is
\emph{ex post}. We identify suspected cover bidders systematically
from participation data, turning a retrospective observation into a
prospective screening tool.

\paragraph{Procurement in developing economies.}
The bid-rigging literature focuses on developed economies, yet
enforcement is hardest where resources are scarce
\citep{sanchezgraells2019screening}. BEC's scale---4.5 million
tender-items, 40 million bids, 11 years, two auction formats---provides a test bed for screens designed to work with minimal data
infrastructure.
